% Build in this directory: pdflatex -interaction=nonstopmode -halt-on-error class6_handout.tex
\documentclass[a4paper,10pt,pdflatex,ja=standard,jafont=ipaex-type1,
  magstyle=real,scale=1]{bxjsarticle}
\usepackage{lmodern}
\usepackage{amsmath,amssymb,mathtools}
\usepackage{geometry}
\geometry{reset,left=22mm,right=22mm,top=20mm,bottom=17mm,
  headheight=13pt,headsep=6mm,footskip=10mm}
\usepackage{xcolor}
\usepackage{enumitem}
\usepackage{array,booktabs,tabularx}
\usepackage{fancyhdr}

\definecolor{ink}{HTML}{193A5C}
\definecolor{accent}{HTML}{237D82}
\definecolor{soft}{HTML}{EAF3F3}
\definecolor{muted}{HTML}{52616B}
\newcommand{\coursename}{解析学 II}
\newcommand{\handoutnumber}{6}
\pagestyle{fancy}
\fancyhf{}
\fancyhead[L]{\small\color{ink}\coursename{}・授業要約と演習}
\fancyhead[R]{\small\color{ink}第\handoutnumber 回}
\fancyfoot[C]{\small\color{ink}\thepage/2}
\renewcommand{\headrulewidth}{0.3pt}
\setlength{\parindent}{0pt}
\setlength{\parskip}{2pt}
\setlength{\abovedisplayskip}{5pt}
\setlength{\belowdisplayskip}{5pt}
\setlength{\abovedisplayshortskip}{3pt}
\setlength{\belowdisplayshortskip}{4pt}
\setlist[itemize]{leftmargin=1.8em,itemsep=1pt,topsep=2pt,parsep=0pt}

\newcommand{\handouttitle}[1]{%
  \begin{center}
    {\Large\color{ink}#1}\par\vspace{3mm}
    {\color{accent}\rule{0.88\linewidth}{0.7pt}}
  \end{center}\vspace{-2mm}}
\newcommand{\handout}[2]{%
  \renewcommand{\handoutnumber}{#1}%
  \handouttitle{第#1回\quad #2}}
\newcommand{\topic}[1]{\par\vspace{5pt}%
  {\large\sffamily\mdseries\color{accent}#1}\par\vspace{2pt}}
\newcommand{\keybox}[1]{\par\vspace{3pt}\begingroup
  \setlength{\fboxsep}{4pt}%
  \colorbox{soft}{\parbox{\dimexpr\linewidth-2\fboxsep\relax}{\small #1}}%
  \endgroup\par\vspace{2pt}}
\newenvironment{problems}
  {\begin{enumerate}[leftmargin=2em,itemsep=3pt,topsep=3pt,parsep=0pt,
    font=\color{accent}]}
  {\end{enumerate}}


\newcommand{\examplelabel}[1]{{\sffamily\mdseries\color{ink}#1}}
\newcommand{\answerroom}{\par\vspace{16mm}}

\begin{document}
\fontsize{10pt}{13.5pt}\selectfont
\setlength{\abovedisplayskip}{4pt}
\setlength{\belowdisplayskip}{4pt}
\handout{6}{重積分の変数変換（ヤコビアン）}
\topic{ヤコビ行列と面積の伸縮率}
座標変換 $T(u,v)=(x(u,v),y(u,v))$ のヤコビ行列と行列式を
$$
 DT=\begin{pmatrix}x_u&x_v\\y_u&y_v\end{pmatrix},\qquad
 J_T=\det DT=x_uy_v-x_vy_u
$$
とする。微小長方形の2辺の像は、$T_u\,du$ と $T_v\,dv$ で近似される。
この平行四辺形の面積は $|J_T|\,du\,dv$ なので
$$
 dx\,dy=|J_T(u,v)|\,du\,dv.
$$
通常の面積では絶対値を取る。$J_T$ の符号は座標の向きの保存・反転を表す。

\topic{変数変換公式と適用条件}
$D,E$ は有界で境界の面積が0、$D=T(E)$ とする。
$T$ は $E$ の近傍で $C^1$ 級、内部で1対1かつ $J_T\ne0$ とする。
$f$ が $D$ 上可積分なら
$$
 \iint_Df(x,y)\,dx\,dy
 =\iint_E f(x(u,v),y(u,v))\,|J_T(u,v)|\,du\,dv.
$$
境界など面積0の部分での重複・退化は積分値に影響しない。
関数・面積要素・積分領域の3つを同時に変換する。

\topic{線形変換の例}
$x=2u+v$、$y=u+3v$ なら $J_T=\det\begin{pmatrix}2&1\\1&3\end{pmatrix}=5$。
$E=[0,1]^2$ の像 $D$ は平行四辺形で
$$
 \operatorname{Area}(D)=\iint_E5\,du\,dv=5,\qquad
 \iint_D(x+y)\,dx\,dy
 =5\int_0^1\int_0^1(3u+4v)\,dv\,du=\frac{35}{2}.
$$

\topic{極座標：円・扇形・円環に適する座標}
$x=r\cos\theta$、$y=r\sin\theta$、$r\ge0$ とすると
$$
 \det\frac{\partial(x,y)}{\partial(r,\theta)}
 =\begin{vmatrix}\cos\theta&-r\sin\theta\\\sin\theta&r\cos\theta\end{vmatrix}=r,
 \qquad dx\,dy=r\,dr\,d\theta.
$$
半径方向の幅 $dr$ と円弧方向の長さ $r\,d\theta$ の積が面積要素になる。
\par\examplelabel{例：半径 $R>0$ の円板 $D_R$。}
$$
 \iint_{D_R}(x^2+y^2)\,dx\,dy
 =\int_0^{2\pi}\int_0^Rr^3\,dr\,d\theta=\frac{\pi R^4}{2}.
$$

\topic{楕円と多対1の変換}
$a,b>0$、$x=ar\cos\theta$、$y=br\sin\theta$ なら $|J_T|=ab r$。
楕円 $x^2/a^2+y^2/b^2\le1$ は $0\le r\le1,\ 0\le\theta<2\pi$ に対応する。
\par
同じ点を複数回覆う変換では、その重複回数も積分に入る。
1対1になる範囲へ制限してから公式を使うのが基本である。
\keybox{ヤコビアンは逆変換のものと取り違えない。極座標では半径の範囲と角度の範囲を図で確認する。}

\newpage
\handouttitle{第6回\quad 演習問題}
計算過程を記し、必要な条件・積分領域・向きを明示すること。
\begin{problems}
\item $T(u,v)=(3u+v,u-2v)$ のヤコビアンを求め、単位正方形の像の面積を計算せよ。行列式の符号の意味も述べよ。\answerroom
\item 極座標を用いて、$\displaystyle\iint_{x^2+y^2\le1}(x^2+y^2)\,dx\,dy$ を求めよ。\answerroom
\item 円環 $D=\{(x,y):1\le x^2+y^2\le4\}$ 上の $\displaystyle\iint_D\frac{dx\,dy}{x^2+y^2}$ を求めよ。\answerroom
\item $a,b>0$、$D=\{(x,y):x^2/a^2+y^2/b^2\le1\}$ とする。変数変換を用いて $D$ の面積と $\iint_Dx^2\,dx\,dy$ を求めよ。\answerroom
\item $u=x+y,\ v=x-y$ とする。$D=\{(x,y):0\le x+y\le2,\ 0\le x-y\le1\}$ 上の $\iint_D(x+y)\,dx\,dy$ を求めよ。\answerroom
\item $T(u,v)=(u^2-v^2,2uv)$ の $|J_T|$ を求めよ。$u^2+v^2\le1$ の像と重複回数を調べ、$\iint_{u^2+v^2\le1}|J_T|\,du\,dv$ が像の面積と一致するか説明せよ。\answerroom
\end{problems}
\end{document}
